\chapter{Setup}
\label{cha:setup}

Ich empfehle im \LaTeX Workshop ein Recipe in den Einstellungen zu hinterlegen um diese Projekt zu Kompilieren:

\begin{lstlisting}[caption=Einstellungen für ein Recipe des LaTeX Workshops]
"latex-workshop.latex.recipes": [
    {
        "name": "pdflatex -> biber -> makeglossaries -> pdflatex",
        "tools": ["pdflatex", "biber", "makeglossaries", "pdflatex"]
    }
],
"latex-workshop.latex.tools": [
    {
        "name": "biber",
        "command": "biber",
        "args": ["%DOCFILE%"]
    },{
        "name": "makeglossaries",
        "command": "makeglossaries",
        "args": ["%DOCFILE%"]
    },{
        "name": "pdflatex",
        "command": "pdflatex",
        "args": [
            "-synctex=1",
            "-interaction=nonstopmode",
            "-file-line-error",
            "%DOC%"
        ]
    }
]
\end{lstlisting}

Alternativ auch über die Command Line:

\begin{lstlisting}[caption=Manuelle ablauf vom Kompilieren]
pdflatex main.tex
biber main
makeglossaries main
pdflatex main.tex
\end{lstlisting}

Wenn Updates der Quellen und des Glossars nicht mehr notwendig sind, kann das Kompilieren auf nur \lstinline|pdflatex main.tex| gekürzt werden.